\section{Метод FieldMathCalculations::ParsePolynomials}
\footnotesize
\begin{listing}{1}
/*************************************************************************
* Метод         : FieldMathCalculations::ParsePolynomials                *
*------------------------------------------------------------------------*
* Назначение    : Парсинг строковых выражений в структуру полиномов      *
* Входные       : a_str, b_str - строки формул полей A и B               *
* Выходные      : Нет (обновляет статические поля класса)                *
*------------------------------------------------------------------------*
* ОПИСАНИЕ ФУНКЦИОНИРОВАНИЯ:                                             *
* Метод преобразует текстовые представления векторных полей              *
* в математическую структуру Polynomial.                                   *
* *
* АЛГОРИТМ:                                                              *
* 1. Вызывает ExtractComponents для каждой строки поля (A и B).          *
* 2. ExtractComponents выполняет нормализацию и разделение на            *
* компоненты X, Y, Z.                                                    *
* 3. Результаты сохраняются в статические переменные класса.             *
*************************************************************************/
void FieldMathCalculations::ParsePolynomials(const wstring& a_str, 
                                             const wstring& b_str) {
    ExtractComponents(a_str, fieldA_X, fieldA_Y, fieldA_Z);
    ExtractComponents(b_str, fieldB_X, fieldB_Y, fieldB_Z);
}
\end{listing}
\normalsize

\pagebreak

\section{Метод FieldMathCalculations::EvaluateFieldA}
\footnotesize
\begin{listing}{1}
/*************************************************************************
* Метод         : FieldMathCalculations::EvaluateFieldA                  *
*------------------------------------------------------------------------*
* Назначение    : Вычисление вектора поля A в заданной точке             *
* Входные       : x, y, z - координаты точки                             *
* Выходные      : Vector3D - результирующий вектор поля в точке          *
*------------------------------------------------------------------------*
* ОПИСАНИЕ ФУНКЦИОНИРОВАНИЯ:                                             *
* Выполняет расчет численных значений компонент векторного поля A        *
* в указанной точке.                                                     *
* *
* АЛГОРИТМ:                                                              *
* 1. Обращается к компонентам fieldA_X, fieldA_Y, fieldA_Z.              *
* 2. Вызывает метод Evaluate для каждого компонента.                     *
* 3. Объединяет результаты в структуру Vector3D.                         *
*************************************************************************/
Vector3D FieldMathCalculations::EvaluateFieldA(double x, double y, double z) {
    return { fieldA_X.Evaluate(x, y, z), 
             fieldA_Y.Evaluate(x, y, z), 
             fieldA_Z.Evaluate(x, y, z) };
}
\end{listing}
\normalsize

\pagebreak

\section{Метод FieldMathCalculations::EvaluateFieldB}
\footnotesize
\begin{listing}{1}
/*************************************************************************
* Метод         : FieldMathCalculations::EvaluateFieldB                  *
*------------------------------------------------------------------------*
* Назначение    : Вычисление вектора поля B в заданной точке             *
* Входные       : x, y, z - координаты точки                             *
* Выходные      : Vector3D - результирующий вектор поля в точке          *
*------------------------------------------------------------------------*
* ОПИСАНИЕ ФУНКЦИОНИРОВАНИЯ:                                             *
* Аналогичен EvaluateFieldA, но работает с полем B, используя методы-    *
* геттеры для доступа к компонентам.                                     *
* *
* АЛГОРИТМ:                                                              *
* 1. Вызывает методы GetFieldBX, GetFieldBY, GetFieldBZ.                 *
* 2. Применяет метод Evaluate для каждого компонента.                    *
* 3. Возвращает вектор в виде структуры Vector3D.                        *
*************************************************************************/
Vector3D FieldMathCalculations::EvaluateFieldB(double x, double y, double z) {
    return { GetFieldBX().Evaluate(x, y, z), 
             GetFieldBY().Evaluate(x, y, z), 
             GetFieldBZ().Evaluate(x, y, z) };
}
\end{listing}
\normalsize

\pagebreak

\section{Метод FieldMathCalculations::Calc\_PartialDerivative}
\footnotesize
\begin{listing}{1}
/*************************************************************************
* Метод         : FieldMathCalculations::Calc_PartialDerivative          *
*------------------------------------------------------------------------*
* Назначение    : Численное вычисление частной производной функции       *
* Входные       : f - функция, variable - индекс переменной, x,y,z       *
* Выходные      : double - приближенное значение производной             *
*------------------------------------------------------------------------*
* ОПИСАНИЕ ФУНКЦИОНИРОВАНИЯ:                                             *
* Вычисляет производную функции f по переменной x, y или z               *
* методом центральной разности.                                          *
* *
* АЛГОРИТМ:                                                              *
* 1. Шаг дифференцирования h = 1e-4.                                     *
* 2. Вычисляет отношение: (f(x+h) - f(x-h)) / (2*h).                     *
*************************************************************************/
double FieldMathCalculations::Calc_PartialDerivative(FUNC3D f, int variable, 
                                                     double x, double y, 
                                                     double z) {
    double h = 1e-4;
    if (variable == 1) return (f(x + h, y, z) - f(x - h, y, z)) / (2 * h);
    if (variable == 2) return (f(x, y + h, z) - f(x, y - h, z)) / (2 * h);
    return (f(x, y, z + h) - f(x, y, z - h)) / (2 * h);
}
\end{listing}
\normalsize

\pagebreak

\section{Метод FieldMathCalculations::IntegrateQuarterDisk}
\footnotesize
\begin{listing}{1}
/*************************************************************************
* Метод         : FieldMathCalculations::IntegrateQuarterDisk            *
*------------------------------------------------------------------------*
* Назначение    : Численное интегрирование функции по четверти диска     *
* Входные       : func - подынтегральная функция, radius, steps          *
* Выходные      : double - значение интеграла                            *
*------------------------------------------------------------------------*
* ОПИСАНИЕ ФУНКЦИОНИРОВАНИЯ:                                             *
* Численное интегрирование двумерной функции по области четверти         *
* диска в полярных координатах (r, phi).                                 *
* *
* АЛГОРИТМ:                                                              *
* 1. Переход к полярным координатам: x = r*cos(phi), y = r*sin(phi). "phi"-   *
*обозначение полярного угла в системе полярных координат. *
* 2. Вычисление двойной суммы по сеткам r и phi с учетом якобиана (r).   *
*************************************************************************/
double FieldMathCalculations::IntegrateQuarterDisk(
    function<double(double, double)> func, double radius, int steps) {
    double dr = radius / steps, dphi = (M_PI / 2.0) / steps, sum = 0.0;
    for (int i = 0; i < steps; i++) {
        double r = (i + 0.5) * dr;
        for (int j = 0; j < steps; j++) {
            sum += func(r * cos((j + 0.5) * dphi), 
                        r * sin((j + 0.5) * dphi)) * r;
        }
    }
    return sum * dr * dphi;
}
\end{listing}
\normalsize

\pagebreak